\chapter{Design}

\section{Architettura}

L'architettura dell'applicazione segue il pattern architetturale \textit{Model--View--Controller}
(MVC), seppur in una sua variante event-driven con \textit{Model} attivo e aggiornamenti
in modalità \textit{push}. Infatti gli aggiornamenti significativi dello stato della partita vengono
comunicati alla scena interessata da parte del Model stesso, mentre per il flusso opposto
gli input dell'utente vengono intercettati dalla View e interpretati dal Controller,
secondo MVC standard.

L'entry-point principale del \textit{Model} è rappresentato dall'interfaccia \texttt{Board\-Game},
che incapsula lo stato e la logica principale della partita. Data la sua composizione ha
il compito di applicare le regole del dominio ed eseguire le conseguenti operazioni sulle altre classi.
L'interfaccia espone metodi adeguati alla modifica del suo stato e alla sua interrogazione, permettendo
al Controller di interagire con esso quando uno dei giocatori desidera applicare una determinata mossa.
È presente tuttavia un secondo punto di accesso, rappresentato dall'interfaccia \texttt{AiTurnService},
che gestisce il servizio di collegamento tra il Controller e l'algoritmo decisionale dell'IA\@.

L'interfaccia \texttt{MainView} funge da entry-point della \textit{View} e rappresenta la finestra principale
dell'applicazione. Essa tiene traccia delle varie scene dell'applicazione ed è in grado di navigare tra loro grazie a un
gestore del layout, mentre il criterio di navigazione è affidato al Controller. Per scena si intende una specifica
schermata del sistema identificata dall'interfaccia \texttt{Scene}. Ciascuna delle scene registrate espone
un'interfaccia provvista di metodi specifici che permettono al Controller di registrare a un determinato
componente grafico una determinata azione. In questo modo la View si occupa della rappresentazione grafica, della
selezione degli elementi e della raccolta degli eventi generati dall'interfaccia, ma non applica direttamente
le regole di gioco. Inoltre \texttt{MatchScene} può essere considerato come un punto d'accesso secondario
della View, in quanto è presente un flusso di comunicazione diretta tra essa e il Controller dedicato allo
stato della partita.

Il \textit{Controller} è rappresentato principalmente dalle interfacce \texttt{MainControl\-ler} e
\texttt{SceneCoordinator}. Il controller coordina il flusso applicativo generale, crea le schermate, registra le scene
nella vista principale e gestisce la navigazione tra di esse, a richiesta delle azioni utente catturate dalla View.
Per evitare un'eccessiva concentrazione di responsabilità, si suddividono le responsabilità in controller
specializzati per le singole schermate. Questi collegano le callback esposte dalle View alla logica
applicativa corrispondente.

\begin{figure}[htbp]
    \centering
    \includegraphics[width=0.80\textwidth]{UMLs/ArchitetturaMVC.pdf}
    \caption{Schema UML dell'architettura MVC dell'applicazione.}
\end{figure}

\section{Design Dettagliato}

\subsection{Massimiliano Ria}

\subsubsection{Adattamento dell'interfaccia a diverse risoluzioni}

\textbf{Problema:} Dopo aver ottenuto un'interfaccia grafica funzionante e
un'architettura coerente, è emersa la necessità di mantenerne l'aspetto
uniforme su risoluzioni e rapporti d'aspetto diversi.

\textbf{\\Soluzione:} Un primo approccio al problema è stato l'introduzione della
classe \texttt{ViewMetrics}, usata come punto di accesso unico alle principali
dimensioni grafiche dell'applicazione. Nella sua versione iniziale, tali misure
venivano calcolate a partire dalla dimensione fisica dello schermo. La classe
determina una sola volta le grandezze necessarie e le espone tramite metodi
statici dedicati, come ad esempio \texttt{cardWidth()} e \texttt{cardHeight()}.
In questo modo si riduce la duplicazione del codice e ogni scena può richiedere
solo la metrica di cui ha bisogno.

I test successivi hanno però evidenziato un limite: scalare ogni componente in
base all'intero schermo non garantiva automaticamente una buona visualizzazione.
Su schermi con risoluzioni elevate o con un rapporto d'aspetto diverso da quello
previsto, alcuni elementi risultavano troppo grandi o disposti in modo poco
bilanciato. Inoltre lo sfondo del menu, realizzato nativamente con risoluzione
3840$\times$2160 e rapporto 16:9, rischiava di essere adattato a superfici non coerenti
con la sua proporzione originale, producendo tagli dell'immagine o distorsioni
impreviste.

Per risolvere il problema è stato introdotto il concetto di \textit{viewport}:
l'applicazione continua ad aprirsi a schermo intero, ma la zona effettiva in cui
vengono disegnate le scene corrisponde alla più grande area 16:9 contenibile
nello schermo dell'utente. Il calcolo è centralizzato nel metodo
\texttt{ViewMetrics.viewportSize()}, che confronta il rapporto tra la
larghezza e l'altezza disponibili con il rapporto obiettivo 16:9.

La classe \texttt{BlackBounds} applica concretamente questa scelta: viene usata
da \texttt{MainViewImpl} come pannello radice della finestra e contiene il
pannello delle scene gestito dal \texttt{CardLayout}. Nell'override del metodo
\texttt{doLayout()}, richiamato automaticamente da Swing quando la finestra
viene mostrata per la prima volta o quando cambia dimensione,
\texttt{BlackBounds} calcola la viewport a partire dalla dimensione corrente
della finestra, centra l'unico componente figlio e gli assegna l'area 16:9
risultante. Tutto lo spazio esterno rimane nero,
generando bande laterali oppure superiori e inferiori in base alla situazione.

\begin{figure}[H]
    \centering
    \includegraphics[width=0.80\textwidth]{UMLs/ViewportAdattamento.pdf}
    \caption{Schema UML dell'adattamento della viewport al rapporto 16:9.}
\end{figure}

\subsubsection{Multi-Risoluzione per schermi Hi-DPI}

\textbf{Problema:} Nelle ultime fasi di sviluppo è emerso un bug legato alla
qualità delle immagini delle carte. Il problema non si presentava in modo
uniforme: su alcuni schermi le texture risultavano nitide, mentre su altri
apparivano sgranate o sfocate, senza che inizialmente fosse riconoscibile un
legame evidente con la risoluzione nominale del monitor.

L'analisi successiva ha ricondotto il comportamento agli schermi Hi-DPI\@. Su
questi dispositivi il sistema operativo può applicare uno \textit{zoom} logico:
l'interfaccia continua a ragionare in pixel logici, ma ogni pixel logico viene
poi rappresentato tramite più pixel fisici. Questa scelta evita che finestre e
componenti diventino troppo piccoli su pannelli molto densi, ma introduce una
differenza tra la dimensione richiesta da Swing e la dimensione realmente usata
dal dispositivo di output. Nel caso specifico, le classi della View ottenevano le
immagini tramite \texttt{ImageLoader}, che le restituiva già scalate alle
dimensioni calcolate da \texttt{ViewMetrics}. Quando però Swing disegnava quelle
icone su uno schermo con fattore di scala maggiore di 1, l'immagine già ridotta
veniva adattata a una superficie fisica più grande, producendo l'effetto di
sgranatura osservato sulle carte.

\textbf{\\Soluzione:} L'idea adottata è stata quella di non trattare più una
richiesta di immagine scalata come una singola bitmap, ma come un'immagine
multi-risoluzione. Per ogni risorsa grafica richiesta a una certa dimensione
logica viene preparata una variante coerente con quella dimensione e, quando
necessario, una seconda variante coerente con la dimensione fisica del
dispositivo. In questo modo le classi della View continuano a chiedere immagini
in termini di pixel logici, ma il sottosistema grafico di Java può selezionare
la variante più adatta al momento del disegno.

Il fix è stato concentrato nel repository usato da \texttt{ImageLoader}. Quando
arriva una richiesta di immagine scalata, \texttt{AsyncCachedImageRepository}
parte dalla dimensione logica richiesta dalla View e ricava anche la dimensione
fisica corrispondente, usando il fattore di scala dello schermo quando
disponibile. Il flusso rimane quindi unico: la View chiede un'immagine di una
certa dimensione in pixel logici; il caricatore prepara anche la versione adatta
ai pixel fisici realmente usati dal monitor.

Da quel punto la risorsa originale viene letta come \texttt{BufferedImage}, così
può essere ridimensionata in modo controllato. Se le due dimensioni coincidono,
viene prodotta una normale icona scalata. Se invece lo schermo Hi-DPI richiede
più dettaglio, vengono create due varianti della stessa immagine: una alla
dimensione logica e una alla dimensione fisica. Queste varianti vengono poi
raccolte in una \texttt{BaseMultiResolutionImage}, incapsulata in una
\texttt{ImageIcon}. In questo modo Java riceve un'unica icona, ma può scegliere
internamente la versione più adatta al dispositivo su cui sta disegnando.

Anche la cache è stata adattata a questo comportamento: le immagini scalate non
sono identificate solo dalla dimensione richiesta dalla View, ma anche dalla
dimensione fisica per cui sono state generate. Questo evita che una bitmap
pensata per uno schermo a scala normale venga riutilizzata su uno schermo
Hi-DPI, o viceversa.

Le classi della View restano quindi semplici e non devono conoscere il dettaglio
della multi-risoluzione. \texttt{CardComponent} e i pannelli che mostrano le
carte continuano a passare da \texttt{ImageLoader} e dalle misure di
\texttt{ViewMetrics}; la differenza è gestita prima del disegno, nel punto in
cui le immagini vengono caricate e precaricate. Rispetto ai refactor valutati,
che spostavano il ridimensionamento più vicino al rendering finale, questa
soluzione è risultata più equilibrata: conserva il precaricamento asincrono già
presente, mantiene invariata l'API della View e delega a Java la scelta della
variante grafica corretta.

\subsubsection{Stato Simulabile per l'IA}

\textbf{Problema:} L'approccio scelto per l'implementazione dell'IA in fase di
progettazione ha da subito evidenziato un concetto inevitabile: per permettere
all'IA di valutare l'applicazione di determinate mosse sullo stato della
partita, era necessario che queste venissero effettivamente applicate sul tavolo
da gioco, ma modificare direttamente lo stato reale del match avrebbe compromesso
la partita in corso.

\textbf{\\Soluzione:} La soluzione più logica è stata dunque introdurre una
rappresentazione separata dello stato della partita, ricostruita in un contesto
simulato e interamente modificabile. Questo ruolo è affidato all'interfaccia
\texttt{AiGameState} e alla sua implementazione \texttt{AiGameStateImpl}, che
espongono le operazioni minime necessarie alla pianificazione: danneggiare un
giocatore, consumare mana, giocare una carta, danneggiare o rimuovere una carta
dal campo e produrre una copia dello stato corrente.

La costruzione dello stato simulabile è delegata a
\texttt{AiGameStateFactory}, implementata da \texttt{AiGameStateFactoryImpl}.
La factory interroga lo stato reale attraverso \texttt{BoardGame} e costruisce
lo snapshot copiando i dati dello stato reale e usando copie delle carte.
Lo stato del giocatore umano viene creato in maniera dedicata per evitare di
copiare dati sensibili che dal punto di vista dell'IA non sono accessibili.

Le informazioni dei singoli giocatori sono isolate dietro \texttt{PlayerState}
e \texttt{PlayerStateImpl}. Questa classe mantiene l'identificativo del
giocatore, la sua salute, il mana attuale, l'eventuale mano e l'armata in campo, mentre
\texttt{CardStateImpl} rappresenta gli stati delle carte copiate.

Infine l'interfaccia \texttt{AiStateTransition} completa l'implementazione, rappresentando
il concetto di transizione da uno stato simulato a un altro.
La classe \texttt{AiStateTransitionImpl} infatti crea una copia dello stato
ricevuto e poi la modifica in base all'azione da applicare:
\texttt{PlayCardAction} consuma mana e sposta la carta dalla mano all'armata,
\texttt{AttackHeroAction} riduce la salute dell'eroe avversario, mentre
\texttt{AttackCardAction} applica i danni a entrambe le creature e rimuove
quelle sconfitte.
Gli algoritmi decisionali possono così avere un riferimento a questa classe e utilizzarla
per esplorare scenari successivi senza dipendere dalle regole concrete di modifica dello stato.

\begin{figure}[H]
    \centering
    \includegraphics[width=0.80\textwidth]{UMLs/AiStatoSimulabile.pdf}
    \caption{Schema UML dello stato simulabile utilizzato dall'intelligenza artificiale.}
\end{figure}

\subsubsection{Politica Decisionale Sostituibile per l'IA}

\textbf{Problema:} A un certo punto dell'implementazione dell'IA, è sorta la necessità
di integrarla con il ciclo reale del turno. L'algoritmo poteva
scegliere una sequenza di mosse su uno stato astratto, ma tali decisioni dovevano
poi essere applicate alla partita effettivamente giocata dall'utente. Inoltre la
selezione della difficoltà richiedeva di poter sostituire la politica
decisionale senza modificare il controller della partita o le regole del modello.

\textbf{\\Soluzione:} La soluzione adottata separa la scelta delle mosse dalla
loro esecuzione reale. Il punto di raccordo principale è rappresentato
dall'interfaccia \texttt{AiTurnService} e dalla classe \texttt{AiTurnServiceImpl}.
Il servizio riceve il \texttt{BoardGame} corrente, usa
\texttt{AiGameStateFactory} per costruire lo snapshot simulabile e delega la
decisione a un oggetto di tipo \texttt{AiDecisionAlgorithm}. Il controller della
partita non deve quindi conoscere il funzionamento interno dell'IA: richiede
soltanto il piano del turno e ottiene una lista ordinata di \texttt{AiAction}.

La sostituibilità della politica decisionale è ottenuta applicando il pattern
\textit{Strategy}. L'interfaccia \texttt{AiDecisionAlgorithm} definisce il
contratto comune degli algoritmi, mentre le implementazioni
\texttt{GreedySequentialAiAlgorithm} e
\texttt{DepthLimitedLookaheadAiAlgorithm} rappresentano due strategie
alternative. La prima sceglie iterativamente la migliore azione che produce un
miglioramento immediato dello stato, la seconda esplora sequenze di azioni fino
a una profondità limitata e seleziona quella con la valutazione euristica
migliore. Entrambe usano gli stessi componenti di supporto:
\texttt{AiActionGeneratorImpl} per generare le azioni legali,
\texttt{AiStateTransitionImpl} per applicarle allo stato simulato e
\texttt{HeuristicAiStateEvaluator} per assegnare un punteggio agli scenari.

La scelta concreta della strategia avviene in \texttt{MainControllerImpl}, dove
una mappa associa \texttt{Difficulty.NORMAL} a
\texttt{GreedySequentialAiAlgorithm} e \texttt{Difficulty.HARD} a
\texttt{DepthLimitedLookaheadAiAlgorithm}. Quando l'utente avvia una partita,
il controller costruisce \texttt{AiTurnServiceImpl} passando l'algoritmo
corrispondente alla difficoltà selezionata. In questo modo l'aggiunta di una
nuova difficoltà o di un nuovo algoritmo richiede soltanto una nuova
implementazione della strategia e la sua registrazione nel punto di creazione.

L'applicazione delle azioni decise dall'IA è invece responsabilità di\\
\texttt{AiAction\-Executor}, implementata da \texttt{AiAction\-Executor\-Impl}. Questa
classe traduce le azioni astratte prodotte dagli algoritmi nelle operazioni
effettive di \texttt{BoardGame}: posizionare una carta, attaccare l'eroe
avversario oppure attaccare una creatura. Per farlo invoca rispettivamente
\texttt{place}, \texttt{attackHero} e \texttt{attackCard}. Il
\texttt{MatchController} rimane così il coordinatore del turno, ma non
contiene logica decisionale né dettagli di conversione tra azioni simulate e
azioni reali.

\begin{figure}[H]
    \centering
    \includegraphics[width=0.80\textwidth]{UMLs/AiPoliticaDecisionale.pdf}
    \caption{Schema UML della politica decisionale sostituibile dell'intelligenza artificiale.}
\end{figure}

\subsection{Simone Marsili}

\subsubsection{Gestione delle carte di gioco}

\textbf{Problema:} Le carte nel gioco possono essere di varie tipologie, come
\texttt{Creatures} e \texttt{Spells} (requisito opzionale del progetto), e il
sistema deve essere estendibile a nuove aggiunte e tipi di carte. Inoltre le
carte devono contenere un apposito identificatore, il quale permette le
operazioni sul modello durante la partita e impedisce la duplicazione di dati
tra istanze. Tuttavia è necessario separare l'identità runtime delle carte
dalla loro definizione statica. Si è cercato inoltre un modo per separare
efficacemente dichiarazioni e implementazioni delle creature.

\textbf{\\Soluzione:} L'interfaccia \texttt{Creature} rappresenta una tipologia
di carte schierabili, e l'implementazione \texttt{CreatureImpl} si compone di
un \texttt{CardId}, il quale informa sulla tipologia di carta e sulla sua
numerazione univoca nel game, e di una \texttt{CreatureDefinition}, che
rappresenta invece l'aspetto statico della carta (le sue statistiche base).
L'implementazione delle creature è stata realizzata tramite una
\textit{Simple Factory}, che rende il codice facilmente estendibile e
manutenibile. La factory utilizza un \texttt{IdGenerator}, che non istanzia
autonomamente, ma che riceve tramite costruttore, riducendo l'accoppiamento e
migliorando testabilità ed estendibilità.

\begin{figure}[H]
    \centering
    \includegraphics[width=0.80\textwidth]{UMLs/Card.pdf}
    \caption{Schema UML della gestione delle carte e delle relative definizioni.}
\end{figure}

\subsubsection{Gestione del database di carte}

\textbf{Problema:} Il gioco prevede la presenza di un database contenente tutte
le carte che possono comparire in un match. Il sistema necessita di una
rappresentazione centralizzata e facilmente estendibile delle carte, caricata
da una fonte esterna (file), mantenendo separata la logica di parsing dalla
logica di dominio. Inoltre, il database deve gestire esclusivamente dati
statici delle carte, risultando indipendente dall'identità runtime delle
istanze generate durante la partita.

\textbf{\\Soluzione:} Il database gestisce esclusivamente definizioni statiche
delle carte (\texttt{CreatureDefinition}), risultando completamente
disaccoppiato dal concetto di identità runtime, che viene invece introdotto
solo in fase di creazione delle istanze.

Inoltre la creazione di un'istanza del database avviene tramite apposita
Factory statica, che utilizza componenti distinti per la lettura del file e per
il parsing del contenuto, mantenendo separate le responsabilità di accesso ai
dati e interpretazione degli stessi.

Le interfacce \texttt{Parser\textless{}T\textgreater{}} e
\texttt{FullReader\textless{}T\textgreater{}} rappresentano due strategie
intercambiabili (pattern \textit{Strategy}) rispettivamente per
l'interpretazione dei dati e per la lettura delle sorgenti, permettendo di
estendere facilmente il sistema a nuovi formati o fonti dati.

\begin{figure}[H]
    \centering
    \includegraphics[width=0.80\textwidth]{UMLs/Database.pdf}
    \caption{Schema UML del database delle carte e dei componenti di lettura e parsing.}
\end{figure}

\subsubsection{Gestione del deck}

\textbf{Problema:} Il gioco richiede la creazione di un mazzo da cui il
giocatore possa pescare progressivamente le carte durante la partita. Tuttavia,
il mazzo non deve essere costruito manualmente né contenere direttamente le
definizioni statiche presenti nel database: ogni carta inserita nel deck deve
essere una nuova istanza runtime, dotata del proprio identificatore.

Un ulteriore problema riguarda la composizione del mazzo. Non tutte le carte
dovrebbero avere necessariamente la stessa probabilità di comparire: ad esempio,
carte più forti possono essere rese meno frequenti rispetto a carte più deboli,
così da ottenere un deck più bilanciato. Il sistema deve quindi permettere di
modificare il criterio di selezione delle carte senza intervenire sulla logica
interna del mazzo o sulla factory che lo costruisce.

\textbf{\\Soluzione:} La gestione del mazzo è stata separata dalla sua
costruzione. L'interfaccia \texttt{Deck} espone soltanto le operazioni
necessarie durante la partita, cioè la pesca di una carta tramite
\texttt{draw()} e la consultazione del numero di carte rimanenti tramite
\texttt{getRemaining()}. L'implementazione \texttt{DeckImpl} incapsula quindi
la lista delle carte effettivamente presenti nel mazzo e si occupa solo della
loro gestione runtime.

La creazione del mazzo è invece affidata a \texttt{DeckFactory}. Questa classe
riceve tramite costruttore un \texttt{CreatureDatabase}, da cui ottiene le
definizioni statiche delle creature, e una \texttt{CreatureImplFactory}, usata
per trasformare ogni \texttt{CreatureDefinition} selezionata in una vera
istanza runtime di \texttt{Card}. In questo modo il database resta separato
dall'identità delle carte usate nella partita, che viene introdotta solo al
momento della generazione del deck.

La classe \texttt{DeckFactory} realizza sostanzialmente una
\textit{Simple Factory}, poiché centralizza la logica di creazione del mazzo e
nasconde al resto del modello i dettagli necessari alla sua costruzione
(recupero definizioni dal database, creazione istanze runtime, selezione
pesata, mescolamento finale\ldots{}).

Il metodo \texttt{createWeighted(size, strategy)} costruisce un mazzo della
dimensione richiesta usando una \texttt{WeightStrategy}. In questa versione
implementata, la factory calcola il peso associato a ciascuna definizione e usa
tale valore per effettuare una selezione casuale non uniforme. In seguito le
carte generate vengono mescolate prima della creazione del \texttt{DeckImpl},
così da evitare che l'ordine del database influenzi l'ordine di pesca.

Nel caso in cui la dimensione richiesta sia maggiore o uguale al numero di
carte presenti nel database, la factory inserisce almeno una copia di ogni
creatura disponibile e completa poi il mazzo tramite selezione pesata. Questo
evita che, in mazzi abbastanza grandi, alcune creature non compaiano mai solo
per effetto della casualità. Quando invece la dimensione richiesta è inferiore
al numero di definizioni disponibili, l'intero mazzo viene generato tramite
selezione pesata.

La soluzione usa il pattern \textit{Strategy} attraverso l'interfaccia
{\texttt{WeightStra\-tegy}}. Tale interfaccia rappresenta l'algoritmo
intercambiabile con cui viene calcolata l'importanza di una
\texttt{CreatureDefinition} durante la costruzione del mazzo.
\texttt{DeckFactory} usa la strategia senza conoscere il criterio concreto
adottato, rendendo possibile introdurre in futuro nuove politiche di
generazione del deck senza modificare la factory. Inoltre la
\texttt{WeightStrategy} è una functional interface, il che rende molto facile
creare al volo delle nuove strategie.

\begin{figure}[H]
    \centering
    \includegraphics[width=0.80\textwidth]{UMLs/Deck.pdf}
    \caption{Schema UML della costruzione e gestione del mazzo di gioco.}
\end{figure}

\subsubsection{Gestione dell'aggiornamento della view}

\textbf{Problema:} Durante lo svolgimento di una partita, lo stato del modello
cambia frequentemente: la partita viene avviata, il turno passa da un giocatore
all'altro, vengono pescate o giocate carte, cambiano i punti vita di creature e
giocatori, varia il mana disponibile e alcune carte possono essere distrutte o
bruciate. Tutti questi cambiamenti devono essere riflessi nella schermata di
gioco in modo coerente e tempestivo.

Una prima soluzione possibile sarebbe stata lasciare al controller il compito di
aggiornare manualmente la view dopo ogni operazione eseguita sul modello. Questa
scelta, però, avrebbe reso il controller responsabile anche dei dettagli di
aggiornamento grafico, aumentando il suo accoppiamento sia con il model sia con
la view.

Un'altra alternativa sarebbe stata usare una notifica generica, basata su un
unico metodo come \texttt{update()}. In questo caso, però, la view avrebbe
ricevuto solo l'informazione che qualcosa nel modello è cambiato, dovendo poi
interrogare il model per ricostruire lo stato aggiornato. Questa variante è
stata scartata perché avrebbe richiesto alla view di conoscere più dettagli del
modello e avrebbe reso meno esplicito il significato dei singoli aggiornamenti.

Il problema era quindi progettare un meccanismo che permettesse al model di
comunicare alla view gli eventi rilevanti della partita, fornendo direttamente
le informazioni necessarie all'aggiornamento grafico, senza costringere la view
a richiedere ogni volta lo stato aggiornato.

\textbf{\\Soluzione:} La soluzione adottata applica il pattern
\textit{Observer} in una variante di tipo \textit{push model}. Il modello della
partita è l'oggetto osservabile, mentre la view della partita è uno degli
osservatori registrati. Quando lo stato del model cambia, il model non invia
una notifica generica, ma comunica direttamente quale evento è avvenuto e quali
dati aggiornati sono necessari per gestirlo.

L'interfaccia \texttt{ObservableGame} definisce il contratto degli oggetti
osservabili, esponendo i metodi per registrare e rimuovere osservatori.

L'interfaccia \texttt{BoardGame} estende \texttt{ObservableGame}, rendendo
esplicito che il modello principale della partita può essere osservato. La
classe {\texttt{BoardGame\-Impl}} implementa questo contratto, mantiene gli
observer registrati e li notifica opportunamente quando si verifica un
cambiamento significativo nello stato della partita.

Gli observer sono rappresentati dall'interfaccia \texttt{GameObserver}. A
differenza di una versione minimale del pattern \textit{Observer}, non è stato
usato un unico metodo generico di aggiornamento. L'interfaccia espone invece
più metodi specifici, ciascuno associato a un evento rilevante del gioco.

Questa scelta rende la comunicazione tra model e view più esplicita. Il nome
del metodo identifica l'evento avvenuto, mentre i parametri trasportano già le
informazioni aggiornate.

Sul lato view, \texttt{MatchView} estende \texttt{GameObserver}. In questo
modo, una schermata di partita è anche un osservatore del modello. La classe
concreta \texttt{MatchScene} implementa \texttt{MatchView} e definisce come
tradurre ciascun evento ricevuto in un aggiornamento grafico.

Il collegamento tra model e view viene effettuato dal controller, che registra
la view come observer del model. Il model, però, non conosce la classe concreta
\texttt{MatchScene}: conosce solo l'interfaccia \texttt{GameObserver}. Questo
mantiene basso l'accoppiamento tra le componenti e rispetta la separazione
prevista dall'architettura MVC\@.

La soluzione ha anche un compromesso: l'interfaccia \texttt{GameObserver} cresce
all'aumentare degli eventi osservabili. Se in futuro venissero introdotti nuovi
eventi di gioco, potrebbe essere necessario aggiungere nuovi metodi
all'interfaccia e aggiornare le classi che la implementano. Nel contesto del
progetto, però, questo costo è accettabile, perché gli eventi rilevanti della
partita sono limitati e ben definiti. In cambio, si ottiene una comunicazione
più leggibile, più controllata e più aderente al dominio del gioco.

\begin{figure}[H]
    \centering
    \includegraphics[width=0.80\textwidth]{UMLs/Observer.pdf}
    \caption{Schema UML del meccanismo di aggiornamento della view tramite observer.}
\end{figure}

\subsection{Francesco Dama}

\subsubsection{Problema: Gestione delle carte possedute e schierate dal giocatore}


\textbf{Problema:} Nel gioco le carte controllate da un giocatore possono trovarsi in stati
differenti durante la partita. Alcune carte sono presenti nella mano del giocatore e rappresentano
le risorse immediatamente giocabili, mentre altre vengono schierate sul campo e diventano parte attiva
del combattimento. Una possibile soluzione sarebbe stata gestire tutte le carte del giocatore tramite
un'unica collezione all'interno della classe PlayerImpl, distinguendo manualmente i diversi stati
delle carte tramite controlli condizionali o flag interni. Questa scelta, però, avrebbe prodotto
un forte intreccio di responsabilità:
\begin{itemize}
\item il giocatore avrebbe dovuto conoscere sia la logica della mano sia quella del board;
\item le operazioni di aggiunta, rimozione e ricerca delle carte sarebbero risultate duplicate;
\item il codice sarebbe diventato più difficile da estendere e mantenere.
\end{itemize}

Inoltre mano e campo di gioco hanno semantiche differenti:
\begin{itemize}
\item la mano rappresenta carte ancora non utilizzate;
\item il board rappresenta creature attive soggette alle regole del combattimento.
\end{itemize}

Il problema era quindi progettare due componenti separate che modellassero chiaramente i diversi
stati delle carte durante la partita, mantenendo però una struttura coerente e facilmente
integrabile nel modello del giocatore.

\textbf{\\Soluzione:} La soluzione adottata separa esplicitamente la gestione delle
carte in due componenti distinte:
\begin{itemize}
\item Hand, responsabile delle carte presenti nella mano del giocatore;
\item Army, responsabile delle creature schierate sul campo.
\end{itemize}

Entrambe le componenti vengono modellate tramite un'interfaccia astratta
\texttt{Hand} e \texttt{Army}
e una rispettiva implementazione concreta \texttt{HandImpl}
e \texttt{ArmyImpl}.

Questa scelta permette di:
\begin{itemize}
\item nascondere i dettagli implementativi;
\item mantenere basso l'accoppiamento;
\item rendere sostituibili le implementazioni concrete;
\item centralizzare la logica di gestione delle collezioni.
\end{itemize}

La classe HandImpl incapsula tutte le operazioni relative alle carte pescate
ma non ancora giocate. Essa si occupa della:
\begin{itemize}
\item memorizzazione delle carte;
\item aggiunta di nuove carte;
\item rimozione delle carte giocate;
\item ricerca tramite identificatore.
\end{itemize}

La gestione della mano è completamente separata dal giocatore, che delega
tali responsabilità alla componente dedicata.
La classe ArmyImpl rappresenta invece le creature attive presenti sul board.
A differenza della mano, l'esercito introduce concetti legati allo stato delle
creature durante il turno, distinguendo ad esempio:
\begin{itemize}
\item creature appena evocate;
\item creature già attive;
\item creature esauste.
\end{itemize}

Questa separazione consente di modellare in modo più aderente il dominio del gioco
e semplifica l'implementazione delle regole di combattimento.

Dal punto di vista implementativo,
entrambe le classi fanno ampio uso di funzionalità moderne di Java:
\begin{itemize}
\item Optional, per rappresentare risultati eventualmente assenti;
\item Stream API, per ricerca e trasformazione delle collezioni;
\item lambda expressions, per definire predicati e operazioni funzionali.
\end{itemize}

Ad esempio, le operazioni di ricerca delle carte vengono effettuate tramite
pipeline funzionali basate su stream, evitando iterazioni imperative
esplicite e migliorando leggibilità e concisione del codice.

La soluzione applica inoltre il principio di composizione:
PlayerImpl non eredita il comportamento di mano o esercito,
ma possiede componenti specializzate dedicate alla loro gestione.

Questo approccio migliora:
\begin{itemize}
\item modularità;
\item estendibilità;
\item separazione delle responsabilità;
\item aderenza all'architettura object-oriented.
\end{itemize}

\begin{figure}[H]
    \centering
    \includegraphics[width=0.7\textwidth]{UMLs/GestioneManoArmy.pdf}
    \caption{Schema UML della gestione di Hand e Army.}
\end{figure}

\subsubsection{Problema: Gestione della pesca delle carte e degli esiti della pesca}

Durante la partita un giocatore può pescare carte dal proprio deck.
Tuttavia l'operazione di pesca non produce sempre lo stesso risultato: il deck potrebbe essere vuoto,
oppure la mano potrebbe essere piena e la carta pescata dovrebbe quindi essere bruciata.
Una possibile soluzione sarebbe stata usare valori speciali (\texttt{null})
oppure lanciare eccezioni per rappresentare i casi particolari della pesca. 
Questa alternativa è stata scartata perché avrebbe reso meno esplicita la semantica 
dell'operazione e avrebbe distribuito la gestione dei casi speciali nel controller o nel model.
Il problema era quindi rappresentare in modo chiaro e tipizzato tutti i possibili esiti dell'operazione
di pesca, mantenendo il codice leggibile ed evitando ambiguità.

\textbf{\\Soluzione:}
La soluzione adottata introduce il tipo \texttt{DrawCardResult}, che rappresenta esplicitamente
il risultato di una pesca. Tale oggetto contiene tutte le informazioni necessarie per descrivere l'esito dell'operazione.
Il tipo di risultato viene modellato tramite l'enumerazione \texttt{DrawCardResultType}, che distingue i vari scenari possibili, come:
\begin{itemize}
\item pesca completata correttamente;
\item deck vuoto;
\item carta bruciata per mano piena.
\end{itemize}

In questo modo il modello comunica chiaramente cosa è accaduto senza ricorrere a valori speciali o controlli impliciti.

La soluzione segue un approccio simile al pattern \textit{Result Object}:
invece di restituire semplicemente una carta o fallire tramite eccezioni,
il metodo restituisce un oggetto che rappresenta semanticamente l'esito dell'operazione.

Questo rende il codice più robusto e leggibile, perché i casi gestibili dal dominio
vengono esplicitati direttamente nel tipo restituito.

\begin{figure}[H]
    \centering
    \includegraphics[width=0.7\textwidth]{UMLs/PoliticaPescaggioMano.pdf}
    \caption{Schema UML della gestione degli esiti della pesca.}
\end{figure}

\subsubsection{Problema:Gestione dello stato del giocatore}

Nel gioco il giocatore rappresenta una componente complessa del modello. Ogni giocatore possiede:
\begin{itemize}
\item un deck;
\item una mano;
\item un esercito di creature schierate;
\item punti vita;
\item mana disponibile;
\item un'identità logica all'interno della partita.
\end{itemize}

Una possibile soluzione sarebbe stata concentrare tutta questa logica direttamente all'interno
della classe del giocatore, rendendola responsabile sia della gestione delle risorse sia della
logica di combattimento e del turno. Questa scelta, però, avrebbe prodotto una classe molto accoppiata e difficile da estendere.

Il problema era quindi modellare il giocatore come aggregatore delle varie componenti del gameplay,
delegando però le responsabilità specifiche a classi dedicate.

\textbf{\\Soluzione:}
La soluzione adottata utilizza l'interfaccia \texttt{Player} come rappresentazione astratta di un giocatore della partita.

La classe \texttt{PlayerImpl} coordina le varie componenti che costituiscono lo stato del giocatore:
\begin{itemize}
\item \texttt{Hand} per la gestione delle carte in mano;
\item \texttt{Army} per le creature schierate;
\item \texttt{Deck} per la pesca;
\item \texttt{ManaManager} per la gestione del mana.
\end{itemize}

Questa scelta permette di suddividere il comportamento in componenti
specializzate, mantenendo il giocatore come coordinatore logico dello stato.
L'identità del giocatore viene separata tramite \texttt{PlayerId}, mentre \texttt{PlayerType}
rappresenta la natura del giocatore (umano o IA). In questo modo il sistema
può distinguere chiaramente il ruolo del giocatore senza introdurre
controlli sparsi nel codice. La creazione dei giocatori viene inoltre
centralizzata tramite \texttt{PlayerFactory}, che costruisce correttamente
tutte le dipendenze necessarie. Questa soluzione riduce l'accoppiamento
e nasconde la complessità di inizializzazione al resto del sistema.

La gestione del mana viene delegata a \texttt{ManaManager}, evitando
che  \texttt{PlayerImpl} debba conoscere i dettagli delle regole di
incremento e consumo del mana durante i turni.

La soluzione applica quindi una combinazione di:
\begin{itemize}
\item composizione;
\item separation of concerns;
\item factory pattern.
\end{itemize}

\begin{figure}[H]
    \centering
    \includegraphics[width=0.95\textwidth]{UMLs/GestioneStatoGiocatore.pdf}
    \caption{Schema UML della gestione dello stato del giocatore.}
\end{figure}

\subsubsection{Problema: Gestione del turno e delle risorse di mana}

Il sistema di turni rappresenta una parte critica della logica di gioco. A ogni cambio turno devono essere aggiornati:
\begin{itemize}
\item il giocatore attivo;
\item il mana disponibile;
\item lo stato di esaurimento delle creature;
\item eventuali operazioni automatiche di inizio turno.
\end{itemize}

Una possibile soluzione sarebbe stata implementare tutta questa logica direttamente all'interno
di \texttt{BoardGameImpl}, centralizzando l'intero flusso della partita in un'unica classe.
Questa scelta, però, avrebbe prodotto una classe troppo grande e difficile da mantenere.

Il problema era quindi isolare la logica specifica di gestione dei turni,
mantenendola separata dalla restante logica del board game.

\textbf{\\Soluzione:}
La soluzione adottata introduce il componente \texttt{TurnManager}, definito internamente a \texttt{BoardGameImpl}, responsabile esclusivamente della gestione dell’alternanza dei turni.

\texttt{TurnManager} coordina:
\begin{itemize}
\item il passaggio del turno tra i giocatori;
\item l'aggiornamento del mana;
\item il reset dello stato delle creature;
\item le operazioni automatiche di inizio turno.
\end{itemize}

Questa separazione permette di mantenere \texttt{BoardGameImpl} focalizzata sulla coordinazione generale
della partita, evitando che conosca tutti i dettagli operativi del turno.

La soluzione migliora:
\begin{itemize}
\item leggibilità;
\item manutenibilità;
\item separazione delle responsabilità.
\end{itemize}

Inoltre rende più semplice modificare in futuro le regole del turno senza impattare il resto del modello.

\begin{figure}[H]
    \centering
    \includegraphics[width=0.80\textwidth]{UMLs/GestioneTurnoRisorse.pdf}
    \caption{Schema UML della gestione del turno del giocatore.}
\end{figure}

\subsection{Lorenzo Mercuriali}

\subsubsection{Coordinamento delle azioni di partita}

\textbf{Problema:} Nel modello di HearthCode una partita coinvolge più
entità di dominio, e il dubbio creatosi nelle prime fasi di progettazione è stato come
rendere coerente il loro aggiornamento e applicare le regole di dominio senza
spargerle in giro per il codice. Le azioni principali richiedevano quindi l'esistenza
di un'entità che contenesse e coordinasse le varie parti create dal sistema.

\textbf{\\Soluzione:} Un primo approccio al problema è stato quello di suddividere queste
responsabilità in due classi differenti, rispettivamente \textit{Game} e \textit{BoardGame}.
La prima avrebbe affrontato interrogazioni di tipo tecnico, come la creazione e l'inizializzazione
della partita, divenendo l'entry point del modello, mentre la seconda avrebbe svolto il ruolo
di wrapper nei confronti delle altre entità del modello. Tuttavia una seconda analisi ha
portato a una conclusione diversa: assegnare a un'entità come \textit{Game} il compito
di stabilire dati su informazioni che concretamente non aveva, obbligando il sistema
a gestire un flusso ridondante, non era affatto sensato. Infatti non rappresentava una responsabilità
diversa, bensì era una facciata finta priva di significato.

La soluzione adottata e poi finalizzata concentra dunque le operazioni di partita
dietro l'unica entità \texttt{BoardGame}, la quale agisce come punto di coordinamento delle
regole di dominio. L'interfaccia è stata modellata attorno alle azioni
che descrivono realmente una partita, così da gestire internamente le operazioni di dominio
ed esporre solo il risultato alle entità richiedenti.

Sul piano implementativo, l'avvio della partita è stato gestito facendo pescare
un numero fisso di carte iniziali, rappresentato dalla costante
\texttt{STARTING\_HAND\_CARDS}. La scelta di mantenere questo valore nella board
permette di legare la regola di apertura del match al punto che coordina lo
stato complessivo, senza duplicarla nelle classi che rappresentano i singoli
giocatori. Anche la verifica della fine della partita resta centralizzata:
\texttt{getWinner()} filtra i giocatori ancora in vita e restituisce un
\texttt{Optional<PlayerId>}, evitando valori speciali o controlli sparsi sul
numero di punti vita.

Le azioni offensive sono state trattate come operazioni atomiche sullo stato di
partita. Nel caso di \texttt{attackCard}, la board recupera l'armata del
giocatore che sta agendo e quella del difensore, verifica che la creatura
bersaglio sia effettivamente presente e che la creatura attaccante possa
attaccare, quindi applica il danno reciproco. Se una delle due creature arriva
a zero punti vita, la rimozione viene eseguita immediatamente sull'armata
corretta. \texttt{attackHero} segue lo stesso principio, ma invece di cercare una
creatura bersaglio applica il valore d'attacco della carta al giocatore
avversario e consuma l'attacco della creatura usata.

Il piazzamento di una carta è stato invece separato in una sequenza di controlli
prima dell'aggiornamento dello stato: la board accetta solo carte di tipo
\texttt{CREATURE} e rifiuta l'azione se l'armata del giocatore corrente è già
piena. Solo dopo queste verifiche la carta può essere rimossa dalla mano del
giocatore e inserita tra le creature in campo. In questo modo le classi
specializzate continuano a gestire i propri dettagli interni, mentre
\texttt{BoardGameImpl} rimane il punto in cui vengono composte le regole che
coinvolgono più entità contemporaneamente.

\subsubsection{Separazione della costruzione della partita tramite Factory}

\textbf{Problema:} Prima che il match possa iniziare devono
essere assemblati componenti diversi del model, e per alleggerire le
responsabilità di \texttt{BoardGame} si è deciso di assegnare il compito della sua
instanziazione al pattern \textit{Factory}.

\textbf{\\Soluzione:} Inizialmente l'assemblaggio era collocato all'interno del costruttore di
\texttt{BoardGameImpl}, rendendo la classe molto verbosa. La board rappresenta
il punto d'ingresso della logica di partita e deve concentrarsi sulla gestione dello stato
e delle azioni di gioco, non sulle scelte di creazione di altre entità.

Una prima alternativa sarebbe stata collocare questa costruzione nel controller
che avvia la partita. In questo modo, il \textit{Controller} avrebbe dovuto incorporare
scelte di costruzione e entità che appartengono al model, perdendo parte del proprio
ruolo di coordinatore e non rispettando il dogma del pattern MVC\@.

La responsabilità di assemblare una partita è stata dunque
isolata in \texttt{BoardGameFactory}, una factory statica riconducibile ai
pattern creazionali e dedicata alla configurazione standard del match. Il
controller non crea direttamente la board né le sue dipendenze: in
\texttt{MainControllerImpl.startMatch()} richiede una \texttt{BoardGame} tramite
\texttt{createDefaultGame()}, ottenendo un oggetto già pronto per essere passato
al \texttt{MatchController}.

La costruzione dei dettagli iniziali è confinata in
\texttt{createDefaultSetup()}. Qui vengono preparati gli oggetti necessari alla
partita: le definizioni delle creature sono caricate tramite
\texttt{CreatureDatabaseFactory}, le carte concrete vengono prodotte con
\texttt{CreatureImplFactory}, i mazzi con \texttt{DeckFactory} e i giocatori
con \texttt{PlayerFactory}. La factory associa poi a ciascun giocatore una
\texttt{ArmyImpl} e stabilisce il giocatore iniziale.

Il risultato dell'assemblaggio non viene passato alla board come una lunga
lista di parametri, ma viene raccolto in \texttt{GameSetup}. Questo oggetto
rappresenta la configurazione iniziale del match e contiene giocatori, armate e
giocatore di partenza, proteggendo la struttura delle collezioni tramite mappe
immutabili. \texttt{BoardGameImpl} riceve un \texttt{GameSetup} nel proprio
costruttore protetto e non decide come costruire i componenti iniziali della
partita, ma usa la configurazione ricevuta per inizializzare il proprio stato
interno, compreso il \texttt{TurnManager}. In questo modo la factory incapsula
la creazione della partita, \texttt{GameSetup} trasporta la configurazione
iniziale e \texttt{BoardGameImpl} rimane focalizzata sulle regole e sulle
operazioni del match.

\begin{figure}[htbp]
    \centering
    \includegraphics[width=0.60\textwidth]{UMLs/BoardGameFactory.pdf}
    \caption{Schema UML della costruzione della partita tramite Factory.}
\end{figure}

\newpage

\subsubsection{Caricamento asincrono e caching delle risorse grafiche}

\textbf{Problema:} Una volta realizzata l'architettura e finalizzata
l'interfaccia grafica dell'applicazione, ci siamo scontrati con un ostacolo significativo. 
La View di HearthCode utilizza molte risorse grafiche, e la loro lettura in diverse parti
del codice ha prodotto dei concreti rallentamenti in alcune schermate del sistema.
Il problema era particolarmente visibile nelle scene più pesanti. Inizialmente
l'ingresso nella partita dopo il click su \textit{Play} poteva congelare l'EDT di
\textit{Swing} per diversi secondi, perché la scena di gioco doveva caricare e ridimensionare
le immagini delle carte nel momento stesso in cui veniva costruita.

Era quindi necessario progettare un meccanismo che permettesse di riutilizzare
le immagini già caricate, avviare in anticipo il caricamento dei gruppi di
risorse associati alle scene principali e lasciare semplice l'uso delle immagini
da parte della View.

\textbf{\\Soluzione:} La soluzione adottata separa l'uso delle immagini dal modo
in cui esse vengono caricate. Il punto di accesso per la View è la facciata
statica \texttt{ImageLoader}, che espone i metodi \texttt{load(String)} e
\texttt{load(String, int, int)} per ottenere rispettivamente immagini originali
e immagini scalate.

Dietro la facciata è presente \texttt{ImageLoaderProxy}, che traduce le
richieste sincrone della View in richieste asincrone verso un
\texttt{ImageRepository}. Il proxy usa il metodo \texttt{join()} solo quando una
classe grafica richiede effettivamente l'icona, ma può anche avviare più
caricamenti in background tramite \texttt{preload()}, combinandoli con
\texttt{Completable\-Future.allOf}. In questo modo il codice della View conserva
un'interfaccia semplice, mentre l'infrastruttura può lavorare in anticipo.

La classe concreta \texttt{AsyncCachedImageRepository} mantiene due cache
concorrenti: una per le immagini originali e una per le immagini scalate.
L'uso di \texttt{computeIfAbsent} fa sì che due richieste uguali condividano lo
stesso \texttt{CompletableFuture}, evitando letture e ridimensionamenti
ripetuti.

Il caricamento effettivo viene eseguito su un \texttt{ExecutorService} dedicato,
creato con thread demoni denominati \texttt{image-loader}. Le immagini originali
sono lette con \texttt{ImageIO}, validate tramite le dimensioni dell'icona e poi
salvate nella cache. Per le richieste scalate, descritte da
\texttt{ImageLoadRequest.scaled}, il repository passa prima dalla cache
dell'immagine originale, caricandola se necessario, e poi produce l'icona alla
dimensione richiesta; le richieste senza ridimensionamento sono invece create con
\texttt{ImageLoadRequest.raw}.

I gruppi di risorse da precaricare sono raccolti in
\texttt{ImagePreloadCatalog}. Il metodo \texttt{menuAndNavigation()}
prepara lo sfondo comune, i banner della schermata finale e le texture dei
pulsanti usati nelle schermate di navigazione; \texttt{match()} prepara le
immagini usate dalle carte nella scena di gioco, compreso il retro delle carte e
le immagini delle creature; \texttt{database()} prepara invece le immagini usate
nella schermata del mazzo. Le dimensioni delle immagini scalate vengono ricavate
da \texttt{ViewMetrics}, così il catalogo produce richieste coerenti con la
risoluzione effettiva dell'interfaccia.

Il precaricamento viene quindi avviato all'avvio dell'applicazione. Qui vengono
chiamati tre punti di ingresso di \texttt{ImageLoader} corrispondenti alle diverse schermate
dell'applicazione e i contenuti grafici da loro richiesti. Ogni metodo
conserva il proprio \texttt{CompletableFuture} statico e lo crea una sola
volta, quindi richieste successive non riavviano lo stesso piano di
caricamento.

In questo modo la View rimane semplice, continuando a richiedere immagini tramite
metodi sincroni; il controller si limita ad avviare i piani di preload;
il repository gestisce caching, esecuzione asincrona e riuso dei risultati.
Quando una risorsa è già presente in cache, la View la ottiene senza ripetere lettura
e ridimensionamento; quando invece il caricamento è ancora in corso, la chiamata sincrona
attende il completamento dello stesso \texttt{CompletableFuture} già avviato.

\begin{figure}[htbp]
    \centering
    \includegraphics[width=\textwidth]{UMLs/CaricamentoRisorseGrafiche.pdf}
    \caption{Schema UML del caricamento asincrono e caching delle risorse grafiche.}
\end{figure}

\newpage
